% =============================================================================
% §30b — Cross-Diagram Node Linking
% =============================================================================
% Specifies the mechanism for drawing edges between nodes that live in
% different cells of a poster/composition.  Superset-only.  Spec-first (2026).
% =============================================================================

\section{Cross-Diagram Node Linking}
\label{sec:cross-diagram-links}

% ─────────────────────────────────────────────────────────────────────────────
\subsection{Concept and Motivation}
% ─────────────────────────────────────────────────────────────────────────────

A poster assembled from independent grammar cells
(Section~\ref{sec:composition}) is, by default, a set of \emph{panels side by
side}: each cell compiles and renders autonomously with no structural
relationship to its neighbours.  Cross-diagram node linking removes that
constraint.  A \texttt{link} statement in the poster DSL draws a directed or
undirected edge between a node in one cell's diagram and a node in another
cell's diagram, turning the poster into a \textbf{linked multi-view system
diagram}.

\begin{description}
  \item[Presentation-layer assertion.]
        A cross-diagram link does not modify either diagram's Domain IR or
        Scene IR.  It is an annotation on the composition itself, resolved
        after every cell has compiled independently.

  \item[Superset-only.]
        The \texttt{link} statement is a new top-level poster keyword with no
        Mermaid equivalent (Section~\ref{sec:superset-extensions}).  A
        standard Mermaid renderer will reject a poster document that contains
        \texttt{link} statements; this is the intended signal that the file
        requires our compiler.

  \item[Additive and non-fatal.]
        Unresolvable links (unknown cell, unknown node, non-linkable diagram
        type) emit a warning and are skipped.  The poster and all valid links
        still render.
\end{description}

% ─────────────────────────────────────────────────────────────────────────────
\subsection{Link Syntax}
\label{sec:cdl-syntax}
% ─────────────────────────────────────────────────────────────────────────────

A \texttt{link} statement is a poster-level declaration written alongside
\texttt{cell} definitions.  Its abstract grammar is:

\begin{lstlisting}[language=,caption={Cross-diagram link statement — abstract grammar}]
link <cellAddr>.<nodeId> <edge> <cellAddr>.<nodeId> [: "label"]

<cellAddr>  ::= "[" row "," col "]"   (* bracket form, 0-indexed *)
             |  colLetters rowNumber  (* Excel form, e.g. A1, B2, AA3 *)

<nodeId>    ::= (* the node identifier as authored in that cell's diagram source *)

<edge>      ::= "-->"                 (* directed arrow *)
             |  "---"                 (* plain undirected *)
             |  "-.-"                 (* dashed undirected *)
             |  "-.."                 (* dashed undirected, alternate *)
             |  "-.->"|"-.->"        (* dashed directed arrow *)
\end{lstlisting}

Cell addresses reuse the dual addressing scheme introduced for the
\texttt{poster} DSL and specified in Section~\ref{sec:composition-posters}:
the bracket form \texttt{[row,col]} is 0-indexed (row first); the Excel form
\texttt{A1}, \texttt{B2}, \ldots{} uses bijective base-26 column letters with
1-indexed row numbers.  Both forms may be mixed freely within a single poster.

The optional \texttt{: "label"} suffix attaches a text label rendered at
the midpoint of the overlay edge, styled by the composition theme's
\texttt{overlayLabelFont} token.

\paragraph{Worked example.}

The following poster specifies an Order Fulfilment system with four cells and
three cross-diagram links that connect nodes across cells:

\begin{lstlisting}[language=,caption={Cross-diagram links in a four-cell poster}]
---
theme: executive
layout: grid 2x2
---
poster "Order Fulfilment System"

  cell A1: flowchart LR
    recv[Receive Order] --> valid[Validate]
    valid --> pay[Payment]
    pay --> pick[Warehouse Pick] --> ship[Dispatch]

  cell B1: classDiagram
    class OrderService
    class PaymentGateway
    class WarehouseSystem
    OrderService --> PaymentGateway
    OrderService --> WarehouseSystem

  cell A2: stateDiagram-v2
    [*] --> Pending
    Pending --> Processing
    Processing --> Shipped
    Shipped --> [*]

  cell B2: sequenceDiagram
    Customer->>OrderService: placeOrder()
    OrderService->>PaymentGateway: charge()
    PaymentGateway-->>OrderService: receipt

link A1.pay --> B1.PaymentGateway : "handled by"
link A1.pick -.-> A2.Processing : "triggers"
link B1.OrderService --- A2.Pending : "creates"
\end{lstlisting}

\texttt{link A1.pay --> B1.PaymentGateway} reads: ``draw a directed arrow
from the node with id \texttt{pay} in cell \texttt{A1} to the class named
\texttt{PaymentGateway} in cell \texttt{B1}''.  The cell addresses
\texttt{A1} and \texttt{B1} are the Excel equivalents of \texttt{[0,0]}
and \texttt{[0,1]}.

% ─────────────────────────────────────────────────────────────────────────────
\subsection{The Node-Anchor Registry}
\label{sec:cdl-anchor-registry}
% ─────────────────────────────────────────────────────────────────────────────

\paragraph{The core engineering prerequisite.}

Grammar layout engines today compute precise node positions internally but
\textbf{discard them}: the output of \texttt{buildFlowScene}, \texttt{buildTreeScene},
and their siblings is a bare \texttt{Scene} containing only flattened
rendering primitives.  For example, the flow layout engine builds a
\texttt{PlacedNode} record for every node (carrying \texttt{x}, \texttt{y},
\texttt{w}, \texttt{h}, \texttt{cx}, \texttt{cy}, and port coordinates) but
returns only the rendered \texttt{Scene}; there is no \texttt{nodeId
$\to$ position} registry in the output type.  \textbf{Making node positions
addressable from the composition layer is the single prerequisite that must be
resolved before any cross-diagram link can be drawn.}

\subsubsection{Candidate Implementations}

Two approaches exist for surfacing node anchors to the composition layer.

\paragraph{Candidate (a) — Populate \texttt{GroupPrimitive.id} and walk the
placed scene.}

The \texttt{GroupPrimitive} in the Scene IR carries an optional \texttt{id?:\
string} field (see \texttt{packages/core/src/scene.ts}).  Each grammar's
layout engine could wrap each node's primitives in a \texttt{GroupPrimitive}
tagged with that node's id.  At composition time the link layer would walk the
placed (already translated-and-scaled) scene, compute the axis-aligned
bounding box of each named group, and use it as the anchor.

\medskip
\textit{Pros:} No change to grammar function return types; \texttt{GroupPrimitive.id}
already exists in the IR.

\medskip
\textit{Cons:}
\begin{itemize}
  \item Recovering a bounding box requires walking the full scene tree for each
        node query (O($n$) primitives per lookup for a scene with $n$ primitives).
  \item The walk must be performed after \texttt{translateAndScale}, meaning
        coordinate reconstruction happens at the composition layer, not the
        grammar layer.
  \item No port information (which side of the node to attach an edge to) is
        recoverable from a bounding box without additional metadata.
  \item Requires every linkable grammar to wrap each node in a named
        \texttt{GroupPrimitive}, which is not currently the case and is a
        fragile convention: a grammar that emits node primitives without a
        wrapping group would silently produce empty anchors.
  \item Not reusable for hit-testing or agent node-addressing without
        additional work.
\end{itemize}

\paragraph{Candidate (b) — Sidecar \texttt{anchors} record on the render
result.}

Each grammar's layout function is extended to return a pair
\texttt{\{ scene: Scene, anchors: NodeAnchorRegistry \}} in addition to the
Scene.  The registry maps each node id to its bounding box and optional
connection ports, expressed in the cell's \textbf{local coordinate space}
(i.e., before composition-level translation and scaling).  Grammars that opt
in populate the registry from data they already hold during layout (the flow
grammar's \texttt{PlacedNode} is a direct source); grammars that have not yet
opted in return an empty registry, and links referencing their nodes degrade
gracefully.

\begin{lstlisting}[language=javascript,caption={NodeAnchorRegistry --- type specification (pseudo-TypeScript)}]
/** Port positions on the four cardinal sides of a node's bounding box. */
type CardinalSide = 'N' | 'S' | 'E' | 'W';

/**
 * Anchor for one node: bounding-box in local cell coordinates,
 * plus optional explicit port attachment points.
 * If ports is absent, the link layer derives port midpoints from
 * (x,y,w,h) at rendering time.
 */
interface NodeAnchor {
  x: number;                                   // bounding-box top-left x
  y: number;                                   // bounding-box top-left y
  w: number;                                   // bounding-box width
  h: number;                                   // bounding-box height
  ports?: Partial<Record<CardinalSide, { x: number; y: number }>>;
}

/** Maps diagram-local node id -> anchor in local cell coordinates. */
type NodeAnchorRegistry = Record<string, NodeAnchor>;

/**
 * Extended render result returned by linkable grammars.
 * Non-linkable grammars continue to return Scene only; the composition
 * layer treats an absent registry as {}.
 */
interface RenderResult {
  scene: Scene;
  anchors: NodeAnchorRegistry;   // empty ({}) for non-linkable grammars
}
\end{lstlisting}

\textit{Pros:} Explicit, typed contract authored by the grammar that knows its
node positions best; no scene walking; carries port information; each grammar
opts in independently; composable independently of the scene.

\medskip
\textit{Cons:} Requires changing grammar function return types; migration is
incremental (one grammar at a time).

\subsubsection{Recommended Implementation: Sidecar Anchors (Candidate b)}

\textbf{The recommended approach is the sidecar} \textbf{\texttt{NodeAnchorRegistry}} \textbf{(candidate b).}  The rationale:

\begin{enumerate}
  \item The grammar layout engine already owns the exact node geometry.  For
        the flow grammar, \texttt{PlacedNode.x}, \texttt{.y}, \texttt{.w},
        \texttt{.h} and the port coordinates \texttt{.rx}/\texttt{.lx}/\texttt{.cx}/\texttt{.by}
        are computed during Phase 4 (coordinate assignment) and available
        before scene emission.  Populating the registry is a two-line loop over
        the placed-node map --- no new computation required.
  \item The composition layer already tracks each cell's
        \texttt{(dx, dy, scale)} from the \texttt{translateAndScale} call in
        \texttt{layoutComposition} (see \texttt{packages/core/src/composition/layout.ts}).
        Transforming local anchors to poster coordinates is exact and trivial:
        \[
          x_{\mathrm{poster}} = x_{\mathrm{local}} \cdot s + d_x, \quad
          y_{\mathrm{poster}} = y_{\mathrm{local}} \cdot s + d_y, \quad
          w_{\mathrm{poster}} = w_{\mathrm{local}} \cdot s, \quad
          h_{\mathrm{poster}} = h_{\mathrm{local}} \cdot s
        \]
        where $s$ is the uniform scale factor and $(d_x, d_y)$ is the
        translation offset already computed by the grid layout pass.
  \item Candidate (a) would require reconstructing the same geometry from the
        scene tree, at the cost of an O($n$)-walk on the flattened primitive
        list, after the coordinate transform has already been applied.
  \item The opt-in migration path means zero risk to existing grammars:
        a grammar that has not yet emitted anchors returns \texttt{\{\}};
        links referencing it degrade to a WARN and are skipped
        (Section~\ref{sec:cdl-semantics}).
\end{enumerate}

\paragraph{Transformation to poster coordinates.}

After each cell's \texttt{RenderResult} is produced, the composition layer
performs the anchor transformation:

\begin{lstlisting}[language=javascript,caption={Anchor transformation to poster coordinates (pseudo-code)}]
function transformAnchors(
  registry: NodeAnchorRegistry,
  dx: number,    // cell's x offset in poster (from translateAndScale)
  dy: number,    // cell's y offset in poster
  scale: number, // uniform scale factor applied to the sub-scene
): NodeAnchorRegistry {
  const out: NodeAnchorRegistry = {};
  for (const [id, anchor] of Object.entries(registry)) {
    const pa: NodeAnchor = {
      x: anchor.x * scale + dx,
      y: anchor.y * scale + dy,
      w: anchor.w * scale,
      h: anchor.h * scale,
    };
    if (anchor.ports) {
      pa.ports = {};
      for (const [side, pt] of Object.entries(anchor.ports)) {
        pa.ports[side as CardinalSide] = {
          x: pt.x * scale + dx,
          y: pt.y * scale + dy,
        };
      }
    }
    out[id] = pa;
  }
  return out;
}
\end{lstlisting}

The resulting poster-coordinate registries from all cells are stored alongside
the placed-cell rectangles for use by the overlay routing pass.

\paragraph{Reusability argument.}

The node-anchor registry is valuable well beyond cross-diagram linking and
should be built regardless of whether the linking feature ships in the first
increment:

\begin{itemize}
  \item \textbf{Hit-testing:} An interactive SVG export can route mouse events
        to the correct diagram node by name, enabling click/hover handlers.
  \item \textbf{Tooltips:} Structured hover labels (from the Domain IR's node
        metadata) can be attached to the exact node bounding box.
  \item \textbf{Agent node-addressing:} The agent integration surface
        (Section~\ref{sec:agent-integration}) currently addresses diagrams by
        their IR structure.  A published anchor registry lets an agent
        reference ``cell A1, node \texttt{pay}'' by logical id rather than by
        pixel coordinates --- a stable, round-trip-safe reference.
  \item \textbf{Accessibility:} ARIA \texttt{role="img"} subtrees can carry
        \texttt{aria-label} attributes derived from each node's anchor id and
        IR label, improving screen-reader navigation of complex diagrams.
\end{itemize}

% ─────────────────────────────────────────────────────────────────────────────
\subsection{Overlay Routing}
\label{sec:cdl-overlay-routing}
% ─────────────────────────────────────────────────────────────────────────────

\paragraph{Overlay layer.}

Cross-diagram edges are rendered as a synthetic pass executed \emph{after} all
cell sub-scenes have been embedded.  The overlay primitives (line/path) are
appended to the poster's primitive list last, placing them above all cell
content in painter's-algorithm order.  No new Scene IR primitive type is
required: overlay edges are \texttt{PathPrimitive} (directed, with arrowhead)
or \texttt{LinePrimitive} (plain) constructed from the same kernel as
intra-diagram edges.

\paragraph{Port selection.}

Given source anchor $A$ (with cardinal ports $A_N$, $A_S$, $A_E$, $A_W$) and
target anchor $B$, the link layer selects the port pair $(p_A, p_B)$ that
minimises the Euclidean distance $\|p_A - p_B\|$.  Tie-breaking order: E
before W before S before N (favouring left-to-right reading direction).  When
explicit ports are absent from the registry, the four cardinal midpoints are
derived from the bounding box:
\[
  A_N = (x + w/2,\; y), \quad
  A_S = (x + w/2,\; y+h), \quad
  A_E = (x+w,\; y+h/2), \quad
  A_W = (x,\; y+h/2).
\]

\paragraph{Phase 1 routing: straight and single-elbow.}

The initial implementation uses two routing strategies selected by the
relative cell positions:

\begin{description}
  \item[Same row, adjacent columns] Straight horizontal segment from $p_A$
        to $p_B$, passing through the inter-cell gap.  This is the most common
        case (left-to-right flow) and produces clean, readable results with no
        routing logic.

  \item[Different rows or non-adjacent columns] Single-elbow (L-shaped) route:
        a horizontal segment from $p_A$ to a midpoint column, then a vertical
        segment to $p_B$'s row, then a horizontal segment to $p_B$.  The
        elbow column is placed at the midpoint between the two cell centres.
        All segments run through the inter-cell gap space and do not overlap
        cell bodies for standard grid layouts.
\end{description}

\paragraph{Phase 2 hardening: orthogonal routing.}

For complex poster layouts where the straight/elbow route would cross a cell
body, the composition engine can activate orthogonal routing:

\begin{enumerate}
  \item Build a set of \emph{forbidden rectangles} from each placed cell's
        bounding box, expanded outward by the theme's \texttt{gap} / 2.
  \item Compute the rectilinear shortest path from $p_A$ to $p_B$ that avoids
        all forbidden rectangles, using the inter-cell gap channels as routing
        channels.
  \item If no route exists (e.g., a cell is surrounded on all sides), emit a
        WARNING and fall back to the Phase 1 route drawn on top of the
        obstructing cell (readable but visually overlapping).
\end{enumerate}

Phase 2 routing is a \emph{hardening step}; Phase 1 covers the vast majority
of practical poster layouts and is the only mode required for the first
increment.

\paragraph{Dimension guard.}

The overlay routing pass applies dimension guards consistent with the
project's overall guard mindset
(Section~\ref{sec:composition-edge-cases}):

\begin{itemize}
  \item If the route must cross more than three cell bodies (unreachable via
        gap channels), emit \texttt{WARNING: link \textit{A.nodeId}
        $\to$ \textit{B.nodeId}: route crosses N cell bodies; overlay may be
        unreadable.}  Still render best-effort.
  \item If both endpoints resolve to the same cell, emit
        \texttt{WARNING: intra-cell link \textit{A.nodeId} $\to$
        \textit{A.nodeId2}: use an intra-diagram edge instead.}  The link is
        skipped (a cross-diagram link within one cell is semantically
        incoherent).
  \item If the computed route length (sum of segment lengths) exceeds three
        times the poster canvas diagonal, emit a WARNING.  This guards against
        degenerate posters where two cells are placed at extreme corners of a
        very large canvas.
\end{itemize}

% ─────────────────────────────────────────────────────────────────────────────
\subsection{Linkable-Type Rules}
\label{sec:cdl-linkable-types}
% ─────────────────────────────────────────────────────────────────────────────

A diagram type is \textbf{linkable} if and only if its layout engine can
produce a non-empty \texttt{NodeAnchorRegistry} --- that is, it authors stable
ids for its visual nodes and can map those ids to bounding boxes.

\begin{table}[h]
\centering\small
\caption{Linkable diagram types (v1)}
\begin{tabularx}{\textwidth}{llX}
\toprule
\textbf{Grammar} & \textbf{Addressable unit} & \textbf{Node id source} \\
\midrule
\texttt{flowchart} / \texttt{graph} & Node & Author-assigned id (e.g.\ \texttt{A}, \texttt{pay}) \\
\texttt{classDiagram} & Class & Class name as written \\
\texttt{stateDiagram-v2} & State & State name / explicit id \\
\texttt{erDiagram} & Entity & Entity name as written \\
\texttt{C4Context} / \texttt{C4Container} / \texttt{C4Component} & Person, System, Container, Component & Alias string \\
\texttt{block-beta} & Block & Block id \\
\texttt{architecture-beta} & Service, Group & Id as written \\
\texttt{mindmap} & Node & Label-derived kebab-slug \\
\texttt{gitGraph} & Commit, Branch & Commit id or branch name \\
\bottomrule
\end{tabularx}
\label{tab:cdl-linkable}
\end{table}

\begin{table}[h]
\centering\small
\caption{Non-linkable diagram types (v1) --- excluded or deferred}
\begin{tabularx}{\textwidth}{llX}
\toprule
\textbf{Grammar} & \textbf{Reason excluded} & \textbf{Degradation} \\
\midrule
\texttt{pie} & No stable node ids; data is label+value pairs & WARN + skip link \\
\texttt{sankey} & Edge-based flow; no persistent node ids & WARN + skip link \\
\texttt{quadrantChart} & Continuous 2-D axes; data points have no ids & WARN + skip link \\
\texttt{radar} & Series labels, not structural nodes & WARN + skip link \\
\texttt{packet-beta} & Byte-field labels, not structural nodes & WARN + skip link \\
\texttt{xychart-beta} & Continuous axis; no node abstraction & WARN + skip link \\
\texttt{timeline} & Events are date-keyed, not node-keyed & WARN + skip link \\
\texttt{sequenceDiagram} & Participants are stable identifiers; deferred to v2 & WARN + skip link \\
\texttt{gantt} & Bars are task-keyed; deferred to v2 & WARN + skip link \\
\bottomrule
\end{tabularx}
\label{tab:cdl-non-linkable}
\end{table}

\paragraph{Rule statement.}

\begin{quote}
A \texttt{link} endpoint is resolvable if and only if (1) the referenced cell
address is valid, (2) the cell's diagram type appears in
Table~\ref{tab:cdl-linkable}, and (3) the specified \texttt{nodeId} is present
in that cell's \texttt{NodeAnchorRegistry} after compilation.  Any endpoint
that fails any condition is \textbf{unresolvable}: the compiler emits one
\texttt{WARNING} per unresolvable endpoint, and the entire \texttt{link}
statement is omitted from the overlay.  A link with one valid endpoint and one
unresolvable endpoint is wholly omitted (a dangling edge has no sensible
routing target).
\end{quote}

% ─────────────────────────────────────────────────────────────────────────────
\subsection{Semantics and Degradation Contract}
\label{sec:cdl-semantics}
% ─────────────────────────────────────────────────────────────────────────────

\begin{quote}
\textbf{Cross-diagram link contract.}
Cross-diagram links are \emph{presentation-layer assertions on the
composition}.  They do not modify, extend, or reference any individual
diagram's Domain IR or Scene IR; each cell compiles as if \texttt{link}
statements did not exist.  Cross-diagram links are \emph{superset-only}: a
standard Mermaid renderer will reject a poster document containing
\texttt{link} statements (the keyword is outside Mermaid's grammar), which is
the intended portability signal.  All link resolution --- cell lookup, node
anchor lookup, port selection, and route computation --- occurs at composition
time, after every cell has compiled independently.  An unresolvable link
(unknown cell address, unknown node id within a valid cell, non-linkable
diagram type, or empty anchor registry) emits \textbf{exactly one
\texttt{WARNING} per unresolvable endpoint} and the link is \emph{silently
omitted}.  The poster renders in full; all other valid links render normally.
\textbf{Cross-diagram links are never fatal.}
\end{quote}

\paragraph{Mermaid portability.}

As with the \texttt{poster} keyword itself (Section~\ref{sec:superset-extensions}),
a file containing \texttt{link} statements cannot be rendered by a standard
Mermaid renderer.  This is intentional and clearly signals the superset
requirement.  Individual cell bodies (the diagram source embedded in each
\texttt{cell} block) remain standard Mermaid syntax and can be extracted and
rendered independently by any Mermaid renderer.

\paragraph{IR integrity.}

Because cross-diagram links are resolved entirely at composition time and
rendered as overlay primitives, they require no changes to any grammar's
Domain IR schema, to the Scene IR primitive types, or to the determinism
contract.  The \texttt{sceneHash} of any individual cell's sub-scene is
byte-identical to what it would be without the \texttt{link} statement.

\paragraph{Agent addressing.}

The \texttt{nodeId} token in a \texttt{link} statement is the same identifier
an agent would use when addressing a node via the structured IR path
(Section~\ref{sec:agent-integration}).  This alignment is intentional: the
anchor registry (Section~\ref{sec:cdl-anchor-registry}) is the shared
mechanism that makes both the cross-diagram link syntax and agent
node-addressing stable and consistent.

% ─────────────────────────────────────────────────────────────────────────────
\subsection{Engineering Prerequisites}
\label{sec:cdl-prereqs}
% ─────────────────────────────────────────────────────────────────────────────

The following work items must be completed before any cross-diagram link can
be drawn.  They are listed in dependency order.

\begin{enumerate}
  \item \textbf{Define \texttt{NodeAnchorRegistry} and \texttt{RenderResult}
        types} in \texttt{packages/core/src/scene.ts} (or a new
        \texttt{scene-anchors.ts}).  This is a pure type addition; no
        behaviour changes.

  \item \textbf{Extend \texttt{buildFlowScene}} to return
        \texttt{RenderResult}.  The flow grammar already computes
        \texttt{PlacedNode} records; emitting the registry requires only
        a loop over the placed-node map at the end of Phase 4.  This is
        the \emph{reference implementation} that verifies the approach
        end-to-end.

  \item \textbf{Extend linkable grammar layout functions} (one per grammar,
        per Table~\ref{tab:cdl-linkable}) to emit anchors.  Each extension is
        independent and can be shipped incrementally.  Non-extended grammars
        return \texttt{anchors: \{\}}.

  \item \textbf{Extend the composition layer} (\texttt{layoutComposition} in
        \texttt{packages/core/src/composition/layout.ts}) to (a) collect anchor
        registries from each cell's \texttt{RenderResult}, (b) apply the
        coordinate transform, and (c) build the poster-coordinate anchor map.

  \item \textbf{Implement the link parser} in the poster DSL
        (\texttt{packages/core/src/frontend/mermaid/poster.ts}): recognise
        \texttt{link} statements and store them alongside the \texttt{cells}
        array in \texttt{PosterDocument}.

  \item \textbf{Implement the overlay routing pass} after the cell embed loop
        in \texttt{layoutComposition}: resolve endpoints, select ports, compute
        routes, and append overlay primitives.

  \item \textbf{Implement graceful degradation}: WARN on unresolvable
        endpoints, skip the link, continue.
\end{enumerate}

Items 1 and 5 are pure additions.  Items 2--4 and 6--7 are incremental
extensions to existing functions with no breaking changes to current callers
(the \texttt{Scene} return type is replaced by \texttt{RenderResult} for
linkable grammars; callers that ignore the \texttt{anchors} field are
unaffected).

% ─────────────────────────────────────────────────────────────────────────────
\subsection*{Cross-references}
% ─────────────────────────────────────────────────────────────────────────────

\begin{itemize}
  \item Composition / Posters (poster DSL, cell addressing):
        Section~\ref{sec:composition} and Section~\ref{sec:composition-posters}
  \item Dual cell addressing (bracket and Excel forms):
        Section~\ref{sec:composition-posters}
  \item Scene IR primitives (\texttt{GroupPrimitive.id},
        \texttt{PathPrimitive}): Section~\ref{sec:scene-ir} and
        \texttt{packages/core/src/scene.ts}
  \item \texttt{translateAndScale} coordinate helper:
        Section~\ref{sec:composition-embed} (paragraph \ref{para:translate-scale-helper})
  \item Superset extension mechanisms (\texttt{poster} keyword,
        superset-only features): Section~\ref{sec:superset-extensions}
  \item Agent integration surface (IR-as-agent-API, node-addressing):
        Section~\ref{sec:agent-integration}
  \item Dimension guard mindset: Section~\ref{sec:composition-edge-cases}
\end{itemize}
